\documentclass{article}

% Official NeurIPS 2026 style, downloaded from the NeurIPS 2026 Call for
% Papers (Formatting_Instructions_For_NeurIPS_2026.zip). No option = the
% anonymised submission format required for double-blind review.
% Use [preprint] for arXiv, [final] for camera-ready.
\usepackage{neurips_2026}

\usepackage[utf8]{inputenc}
\usepackage[T1]{fontenc}
\usepackage{hyperref}
\usepackage{url}
\usepackage{booktabs}
\usepackage{amsfonts}
\usepackage{amsmath}
\usepackage{amssymb}
\usepackage{nicefrac}
\usepackage{microtype}
\usepackage{xcolor}
\usepackage{graphicx}

% GENERATED by paper/gen_numbers.py from the claims registry.
% Do not hand-edit. Re-run the generator instead.

\newcommand{\NrowsRaw}{850}
\newcommand{\NrowsDropped}{33}
\newcommand{\MinAccBefore}{20}
\newcommand{\Nrows}{817}
\newcommand{\Nckpt}{38}
\newcommand{\Nfam}{8}
\newcommand{\Nsch}{6}
\newcommand{\MinCkpt}{3}
\newcommand{\NpredNums}{18}
\newcommand{\AdvWorst}{-39.5}
\newcommand{\AdvBadRows}{36}
\newcommand{\AdvBadOver}{16}
\newcommand{\AdvCtrlRows}{9}
\newcommand{\AdvCtrlOver}{2}
\newcommand{\AdvMaxB}{1.5}
\newcommand{\CtrlRatio}{2.6}
\newcommand{\CtrlGap}{-3.78}
\newcommand{\CalibRatio}{171}
\newcommand{\CtrlZmin}{2.9}
\newcommand{\CtrlZmax}{5.3}
\newcommand{\BaseQwenSmall}{47.20}
\newcommand{\PubQwenSmall}{47.42}
\newcommand{\BaseQwenMid}{59.57}
\newcommand{\PubQwenMid}{60.98}
\newcommand{\ProtoSens}{5.70}
\newcommand{\ExcessN}{6}
\newcommand{\IntervalsExclZero}{0}
\newcommand{\BigLossBelow}{19}
\newcommand{\BigLoss}{25}
\newcommand{\CellsBlocked}{4}
\newcommand{\PooledOneSided}{95.1}
\newcommand{\PubOthersVerifiableCards}{1}
\newcommand{\ProspInside}{118}
\newcommand{\ProspN}{131}
\newcommand{\ProspCov}{90.1}
\newcommand{\ProspCILo}{83.6}
\newcommand{\ProspCIHi}{94.6}
\newcommand{\MaeGain}{0.03}
\newcommand{\TargetMean}{-0.39}
\newcommand{\TargetSd}{1.23}
\newcommand{\NoiseNBench}{11}
\newcommand{\NoiseNThird}{8}
\newcommand{\NoiseMinPct}{15}
\newcommand{\GsmNoiseSd}{2.1}
\newcommand{\BandEmpCov}{93.4}
\newcommand{\BandConfCov}{89.6}
\newcommand{\InfShare}{43}
\newcommand{\PubIntelPaired}{9}
\newcommand{\NCellsInsuff}{8}
\newcommand{\NCellsTotal}{17}
\newcommand{\NCellsRefused}{0}
\newcommand{\ProspSchemes}{4}
\newcommand{\ProspAllRows}{186}
\newcommand{\ProspAllCov}{90.3}
\newcommand{\ProspClusLo}{83.2}
\newcommand{\ProspClusHi}{94.7}
\newcommand{\ProspClusN}{19}
\newcommand{\ProspRowsWfour}{80}
\newcommand{\ProspCkptsWfour}{11}
\newcommand{\ProspCovWfour}{87.5}
\newcommand{\ProspRowsWeightint}{36}
\newcommand{\ProspCkptsWeightint}{5}
\newcommand{\ProspCovWeightint}{91.7}
\newcommand{\ProspRowsFpdyn}{11}
\newcommand{\ProspCkptsFpdyn}{2}
\newcommand{\ProspCovFpdyn}{100.0}
\newcommand{\ProspRowsWeightsixteen}{4}
\newcommand{\ProspCkptsWeightsixteen}{1}
\newcommand{\ProspCovWeightsixteen}{100.0}
\newcommand{\MaeGainLo}{0.01}
\newcommand{\MaeGainHi}{0.05}
\newcommand{\MaeGainFamsPos}{7}
\newcommand{\NoiseMaxPct}{124}
\newcommand{\NoiseNOver}{1}
\newcommand{\NoiseNSmall}{5}
\newcommand{\RefuseBelow}{85}
\newcommand{\BootLoWsixteenBig}{47}
\newcommand{\BootHiWsixteenBig}{100}
\newcommand{\PubChecked}{10}
\newcommand{\PubCards}{260}
\newcommand{\PubOthersCards}{222}
\newcommand{\PubOthersCount}{9}
\newcommand{\PubOthersPaired}{10}
\newcommand{\PubOthersVerified}{6}
\newcommand{\PubCtrlVerified}{147}
\newcommand{\LoFpeight}{-1.11}
\newcommand{\HiFpeight}{+0.48}
\newcommand{\WorstFpeight}{-1.90}
\newcommand{\SevFpeight}{0.0}
\newcommand{\NFpeight}{80}
\newcommand{\LoWeightint}{-1.71}
\newcommand{\HiWeightint}{+1.07}
\newcommand{\WorstWeightint}{-3.32}
\newcommand{\SevWeightint}{1.0}
\newcommand{\NWeightint}{194}
\newcommand{\LoFpdyn}{-1.44}
\newcommand{\HiFpdyn}{+1.34}
\newcommand{\WorstFpdyn}{-8.72}
\newcommand{\SevFpdyn}{1.0}
\newcommand{\NFpdyn}{200}
\newcommand{\LoWeightsixteen}{-0.98}
\newcommand{\HiWeightsixteen}{+0.81}
\newcommand{\WorstWeightsixteen}{-3.04}
\newcommand{\SevWeightsixteen}{1.2}
\newcommand{\NWeightsixteen}{86}
\newcommand{\LoWfour}{-2.87}
\newcommand{\HiWfour}{+1.40}
\newcommand{\WorstWfour}{-8.86}
\newcommand{\SevWfour}{5.7}
\newcommand{\NWfour}{193}
\newcommand{\LoNvfp}{-4.10}
\newcommand{\HiNvfp}{+1.87}
\newcommand{\WorstNvfp}{-8.12}
\newcommand{\SevNvfp}{15.6}
\newcommand{\NNvfp}{64}
\newcommand{\MaeSchememean}{0.7227}
\newcommand{\MaeGlobalmean}{0.7545}
\newcommand{\MaeSchemebench}{0.7607}
\newcommand{\MaeRidge}{0.7639}
\newcommand{\MaeBenchmean}{0.7647}
\newcommand{\MaeGradboost}{0.7878}
\newcommand{\CovCellWfourSmall}{68.8}
\newcommand{\OneSidedCellWfourSmall}{70.1}
\newcommand{\CkptsCellWfourSmall}{2}
\newcommand{\RowsCellWfourSmall}{77}
\newcommand{\CovCellWsixteenBig}{73.7}
\newcommand{\OneSidedCellWsixteenBig}{81.1}
\newcommand{\CkptsCellWsixteenBig}{5}
\newcommand{\RowsCellWsixteenBig}{217}


\title{Calibrated Risk Envelopes for LLM Quantization,\\
and the Limits of Predicting Compression Damage}

% Anonymised for double-blind review. The NeurIPS style renders
% "Anonymous Author(s)" automatically when the [final] option is absent.
\author{%
  Anonymous Author(s) \\
  Affiliation \\
  Address \\
  \texttt{email} \\
}

\begin{document}
\maketitle

\begin{abstract}
\noindent
Quantizing a large language model reduces serving cost and costs some accuracy.
Determining how much normally requires re-running a full benchmark suite, which
is the expense quantization was meant to avoid. We ask whether the accuracy
delta can be predicted instead. Using \Nrows{} published evaluations spanning
\Nckpt{} base checkpoints, \Nfam{} model families and \Nsch{} quantization
schemes, we report three results. First, a \emph{negative} one: per-model
prediction carries almost no signal. Under leave-one-family-out validation the
best learned predictor beats a global mean by \(0.03\)pp of MAE, and both ridge
regression and gradient boosting are \emph{worse} than guessing the average,
because 33--100\% of the variance on most benchmarks is evaluation noise.
Second, a positive one: conditioning only on the quantization scheme yields
split-conformal intervals with \(90.1\%\) empirical coverage
(\(118/131\), 95\% CI \([83.6,94.6]\)) on checkpoints never used in
construction. Third, and most consequential for practice: the published corpus
is outcome-truncated. Every configuration in it is one that someone chose to
release. We rented an A100 and measured deliberately-faulted configurations,
observing losses to \(\AdvWorst{}\)pp against a published worst case of
\(\WorstWfour{}\)pp---measured, we stress, only on models of \AdvMaxB{}B and
below, and on configurations engineered to fail rather than typical ones. We argue the resulting artifact should be described as a
calibrated historical baseline rather than a predictor, and we state precisely
why its coverage guarantee does not transfer to a practitioner's own untuned
recipe. Code, data and every figure's provenance are public.
\end{abstract}

\section{Introduction}

Post-training quantization is now routine. A practitioner choosing between
W4A16, W8A8 and FP8 wants one number: how much accuracy will this cost? The
honest answer today is that nobody knows without running the evaluation.

This paper asks whether that number can be forecast from the public record.
Model publishers release quantized checkpoints with benchmark tables comparing
the compressed model against its dense parent. Aggregated, these are thousands
of paired measurements of exactly the quantity of interest.

Our contributions:

\begin{enumerate}
  \item \textbf{A negative result} (\S\ref{sec:negative}). Per-model prediction
    of the accuracy delta does not work, and we show why: the target is
    dominated by evaluation noise. We report this as the primary finding rather
    than burying it.
  \item \textbf{A calibrated envelope that does work} (\S\ref{sec:results}),
    conditioned on scheme alone, with measured coverage on held-out families
    and an explicit refusal rule for cells where coverage is poor.
  \item \textbf{A measurement of the corpus's own selection bias}
    (\S\ref{sec:tail}), via GPU experiments on deliberately-faulted
    configurations---data the published record structurally cannot contain.
  \item \textbf{An audit} (\S\ref{sec:audit}) recording every error we found in
    our own work, including two that changed published claims.
\end{enumerate}

\section{Related Work}

\paragraph{Predicting evaluations.} BenchPress~\cite{zeng2026benchpress}
establishes the mechanism closest to ours: predict an eval result, estimate
whether to trust it, wrap it in a conformal interval, and decide whether to skip
the real run. Their score matrix over 84 frontier models and 133 benchmarks is
approximately rank-2. We do not claim novelty against them on mechanism. Their
estimand is a model's absolute score; ours is the paired delta between a model
and its compressed self, and their corpus contains no quantized checkpoints.

\paragraph{Conformal prediction on compressed models.}
Tong et al.~\cite{tong2026compression} apply conformal prediction directly to
quantized and sparse LLMs, which is a closer domain than BenchPress. The
estimand differs: their conformal sets are over label space, constructed from a
compressed model's own output probabilities, so the method measures uncertainty
\emph{after} running the model. Ours is constructed over historical accuracy
deltas to give an interval \emph{before} running anything. Their Figure~2 plots
the same quantity we predict, \(\text{Acc}_{\text{compressed}} -
\text{Acc}_{\text{dense}}\), but measures it. Their finding that compression
decouples accuracy from uncertainty bounds what an accuracy-only estimate like
ours can claim (\S\ref{sec:limits}).

\paragraph{Quantization methods.} The checkpoints we harvest are produced with
GPTQ~\cite{frantar2022gptq} and SmoothQuant~\cite{xiao2022smoothquant} via
llm-compressor~\cite{llmcompressor}.

\paragraph{The source of our data.} Kurti\'c et al.~\cite{kurtic2024halfmillion}
report over 500{,}000 evaluations of quantized Llama-3.1 models, finding
recovery above 99\% of baseline. That work does not address whether the
published set of quantized models is a biased sample. Characterising that gap is
the contribution of \S\ref{sec:tail}.

\paragraph{Selection bias.} Survivorship bias in benchmark corpora has been
formalised and quantified in information retrieval~\cite{gupta2022survivorship}.
Conformal prediction's guarantee is known to fail under selection:
conditional on a selection event, calibration data are no longer exchangeable
with a test point~\cite{barber2023beyond,jin2024focal}. We apply this directly
to our own artifact in \S\ref{sec:limits}.

\section{Data}\label{sec:data}

We harvest published model cards into rows of
\((\text{model}, \text{scheme}, \text{benchmark},
\text{acc}_{\text{before}}, \text{acc}_{\text{after}})\), yielding \Nrows{}
rows after filtering, across \Nckpt{} base checkpoints, \Nfam{} families and
\Nsch{} schemes.

\paragraph{An integrity gate.} Cards typically print a \emph{Recovery \%}
column alongside the before/after pair. This is redundant information, and we
exploit it: for every three-column row we check
\(\text{recovery} \approx 100 \times \text{after}/\text{before}\) to a tolerance
derived from the card's own printed precision, and reject rows that disagree.
This caught fabricated values arising from column-order variation between
publishers, and internally inconsistent cards. Rejections are recorded rather
than silently dropped.

\paragraph{A structural limitation, stated at the outset.} All rows come from a
single publisher using a single toolchain. Holding out a model family tests
transfer across \emph{models}; nothing in this design tests transfer across
\emph{quantization practitioners}. \S\ref{sec:tail} gives direct evidence that
this distinction matters.

\section{Method}\label{sec:method}

\paragraph{Split conformal.} Given calibration residuals
\(\{s_i\}_{i=1}^{n}\) and level \(\alpha\), the interval half-width is the
\(\lceil (n+1)(1-\alpha) \rceil\)-th smallest score, which gives exact
finite-sample coverage under exchangeability. Where \(n\) is too small for that
index to exist, we return \(\infty\) rather than a misleadingly finite band.

\paragraph{Mondrian (group-conditional) calibration.} A single marginal
interval advertised at 90\% was in fact a 68\% interval for NVFP4 users---the
users who most need it. We therefore calibrate per scheme.

\paragraph{A refusal rule.} For each (scheme, size-band) cell we measure
coverage directly. Where measured coverage falls below a threshold on a
sufficient number of scored rows, the tool emits
\texttt{INSUFFICIENT CALIBRATION} and no interval. Two cells are refused:
\texttt{w4a16|<2B} at \CovCellWfourSmall\% over \RowsCellWfourSmall{} rows, and
\texttt{w8a16|>10B} at \CovCellWsixteenBig\% over \RowsCellWsixteenBig{} rows.

\paragraph{A support floor.} Coverage measured on many benchmarks from one
checkpoint is not independent evidence; it is one quantization run viewed many
ways. Cells resting on fewer than \MinCkpt{} distinct checkpoints cannot reach
the top tier regardless of measured coverage. \CellsBlocked{} cells are affected.

\section{The negative result}\label{sec:negative}

Under leave-one-family-out validation, row-weighted MAE:

\begin{center}
\begin{tabular}{lcc}
\toprule
predictor & MAE (pp) & beats global mean? \\
\midrule
per-scheme mean (shipped) & \MaeSchememean & yes \\
global mean (baseline)    & \MaeGlobalmean & --- \\
per-(scheme \(\times\) benchmark) & \MaeSchemebench & \textbf{no} \\
ridge, 38 features        & \MaeRidge      & \textbf{no} \\
per-benchmark mean        & \MaeBenchmean  & \textbf{no} \\
gradient boosting         & \MaeGradboost  & \textbf{no} \\
\bottomrule
\end{tabular}
\end{center}

Ridge and gradient boosting are both \emph{worse} than predicting the global
average. Adding model size, bit-width, method, benchmark identity and headroom
all degraded out-of-family performance.

\paragraph{Why.} Estimating per-benchmark evaluation noise from near-lossless
schemes---whose true delta should be approximately zero---gives a noise floor
that accounts for 33--100\% of observed variance on most benchmarks. On GSM8K
the noise alone is \(\pm 2\)pp, while the quantity being predicted has mean
\(-0.36\)pp. The target is a difference of two noisy quantities that largely
cancel: the numerically smallest and hardest-to-predict summary of quantization
damage available.

\paragraph{Implication.} The entire fitted artifact is \NpredNums{} numbers.
A raw per-scheme historical quantile band, with no model and no conformal
machinery, achieves comparable coverage. What conformal contributes is a
finite-sample guarantee and a principled rule for small-\(n\) schemes, not
accuracy. We therefore describe the result as a \emph{calibrated historical
baseline}, not a predictor.

\section{Results}\label{sec:results}

\begin{center}
\begin{tabular}{lcccc}
\toprule
scheme & 90\% interval (pp) & worst observed & \% losing \(>\)3pp & evals \\
\midrule
fp8         & [\LoFpeight, \HiFpeight]           & \WorstFpeight{}pp      & \SevFpeight\%      & \NFpeight \\
w8a8\_int   & [\LoWeightint, \HiWeightint]       & \WorstWeightint{}pp    & \SevWeightint\%    & \NWeightint \\
fp8\_dynamic& [\LoFpdyn, \HiFpdyn]               & \WorstFpdyn{}pp        & \SevFpdyn\%        & \NFpdyn \\
w8a16       & [\LoWeightsixteen, \HiWeightsixteen]& \WorstWeightsixteen{}pp& \SevWeightsixteen\%& \NWeightsixteen \\
w4a16       & [\LoWfour, \HiWfour]               & \WorstWfour{}pp        & \SevWfour\%        & \NWfour \\
nvfp4       & [\LoNvfp, \HiNvfp]                 & \WorstNvfp{}pp         & \SevNvfp\%         & \NNvfp \\
\bottomrule
\end{tabular}
\end{center}

\paragraph{Prospective validation.} After freezing the method we harvested
checkpoints not used during development. Under a strict protocol that excludes
any family present in training and any card whose format informed a parser fix,
coverage is \(118/131 = 90.1\%\), 95\% Clopper--Pearson CI \([83.6, 94.6]\),
against a nominal 90\%. A from-scratch, standard-library-only reimplementation
sharing no code with the pipeline reproduces this figure with zero
disagreements.

\paragraph{What the envelope cannot do.} \IntervalsExclZero{} of \Nsch{}
intervals exclude zero: the tool can never state that a scheme \emph{will} cost
accuracy. And \BigLossBelow{} of \BigLoss{} observed losses worse than 3pp fell
\emph{below} their interval floor---when real damage occurs, the interval
usually misses it. This is the single most important caveat in the paper.

\section{The measured left tail}\label{sec:tail}

The corpus contains only configurations someone chose to publish. To measure
what it structurally cannot show, we rented an A100 and quantized real models
under deliberately-faulted recipes.

\paragraph{Result.} Worst measured loss \(\AdvWorst{}\)pp, against
\(\WorstWfour{}\)pp as the worst loss anywhere in the published corpus.
\AdvBadOver{} of \AdvBadRows{} faulted rows lost 3pp or more; separately,
\AdvCtrlOver{} of \AdvCtrlRows{} rows from our correctly-configured control arm
did too. The dominant single effect was wrong scale granularity---per-tensor
instead of per-group.

\paragraph{Scope, stated plainly.} This tail was measured only on models of
\AdvMaxB{}B and below. It establishes that catastrophic damage is reachable; it
is not a measured bound for larger models.

\paragraph{A confound we found and did not resolve.} Re-running MMLU under the
published protocol reproduces published base accuracy (\BaseQwenSmall{} against
\PubQwenSmall{}; \BaseQwenMid{} against \PubQwenMid{}), closing the harness
question. It exposed a larger one: our own \emph{correctly-configured} control
is \CtrlRatio\(\times\) more damaging than published W4A16 on closely comparable
checkpoints (mean gap \CtrlGap{}pp, \(z = \CtrlZmin\) and \(\CtrlZmax\)). The
cause is approximately \CalibRatio\(\times\) less calibration data than
production defaults, not a bug. Consequently, excess-over-control measures
\emph{bad recipe versus mediocre recipe}, not versus production. It is a
conservative anchor, not a bound, and we decline to publish a corrected point
estimate: the correction moves by \ProtoSens{}pp depending on the anchor set,
and we have only \ExcessN{} rows from two models.

\section{Limitations}\label{sec:limits}

\paragraph{The guarantee is conditional, and void where it matters most.}
Split conformal gives exact coverage under exchangeability between calibration
and test. Our calibration set is the set of configurations a publisher chose to
release---a selection event. Conditional on selection, calibration points are
not exchangeable with an arbitrary test point and the guarantee does not
transfer~\cite{barber2023beyond,jin2024focal}. It holds for a new checkpoint
drawn from the same publication process. It is \textbf{void} for a recipe the
reader tuned themselves and that nobody would have published---which is the case
a practitioner most wants covered.

\paragraph{Accuracy is not the whole of deployment risk.}
Tong et al.~\cite{tong2026compression} show compression frequently decouples
accuracy from uncertainty: a model can hold accuracy while its prediction sets
inflate. A scheme this work calls low-risk on accuracy may still degrade
uncertainty quality. We measure none of that.

\paragraph{One publisher, one toolchain, and why.} See \S\ref{sec:data}. Our own
control arm (\S\ref{sec:tail}) is direct evidence that practitioner and
calibration volume can dominate scheme choice. We measured how confining this
is: surveying \PubChecked{} publishers and \PubCards{} quantized model cards,
\PubOthersPaired{} of \PubOthersCards{} cards from the other
\PubOthersCount{} publishers carry paired before/after evaluations, so others
do publish them. What no other publisher in the sample provides is a
\emph{Recovery} column, which is what lets each row be checked against
\(100 \times \text{after}/\text{before}\) and rejected on disagreement: of
their rows, \PubOthersVerified{} pass that gate, against \PubCtrlVerified{}
from the control publisher. The corpus is single-publisher not because others
publish nothing, but because only one publishes in a self-verifying form.

\paragraph{Rows are not independent.} Up to sixteen benchmark rows come from a
single quantization run. We mitigate with group-conditional calibration and the
support floor, and note that a cluster bootstrap over checkpoints does not
materially widen the headline interval, but the unit of analysis remains a row.

\section{Audit}\label{sec:audit}

We record errors found in our own work, including those that changed published
claims. This section exists because the alternative---reporting only what
survived---would misrepresent how the result was obtained.

\paragraph{Errors that changed a published number.}
\begin{itemize}
  \item \emph{A retracted headline.} An early claim that held-out models came
    from families never used in construction was wrong; they were new
    checkpoints of families already present. The strict figure is \(118/131\).
  \item \emph{Two parser bugs.} A benchmark label typo cost 11 training rows;
    separately, \texttt{MATH-500} was matched as \texttt{Math-Lvl-5},
    contaminating 7 test rows with scores from a different benchmark.
  \item \emph{A centring mismatch.} Calibration measured residuals about the
    stratum mean while intervals centred on the scheme mean, mis-sizing every
    calibrated width. Found only by independent reimplementation.
  \item \emph{A column-order fault.} One publisher orders columns
    \([\text{recovery}, \text{base}, \text{quant}]\), producing a fabricated
    \(-68\)pp delta. The integrity gate rejected 10 of 11 rows on that card,
    passing only the one where \(\text{base} = \text{quant}\) made the check
    degenerate.
\end{itemize}

\paragraph{A regression we nearly shipped.} An audit indicated that a raw
asymmetric empirical band covered better than the shipped symmetric conformal
interval (93.4\% versus 89.6\%). Since quantization damage is left-skewed, a
symmetric band is obviously the wrong shape, and we were about to switch. On
measuring first, the empirical band returns an \emph{infinite} interval on 43\%
of rows, because a two-sided empirical index requires more calibration points
than are often available. An infinite interval covers by construction. On rows
where both are finite, the original wins on coverage \emph{and} width. The
change was reversed.

\paragraph{An overclaimed guarantee.} The artifact was described as carrying a
coverage guarantee without qualification. Checking the method against the
conformal literature rather than against our own earlier work showed the
guarantee is conditional on exchangeability that selection breaks. Corrected in
\S\ref{sec:limits}.

\paragraph{Provenance discipline.} Two parser fixes were informed by inspecting
rows that were already in the test set. Rather than silently retain them, we
record which fixes were test-informed and exclude the affected family from the
strict headline.

\paragraph{Mechanical verification.} Every figure in this paper is a macro
generated from a claims registry that recomputes each value from source
artifacts; no number here is hand-typed. The accompanying repository enforces
this with a test that fails if any outward-facing document contains an untraced
numeric literal.

\section{Conclusion}

Per-model prediction of quantization accuracy damage does not work on published
data, and the reason is measurement noise rather than model capacity. What does
work is narrower: a per-scheme calibrated envelope, validated at
\(90.1\%\) coverage on unseen checkpoints, that refuses to answer where its own
measured coverage is poor. The more useful contribution may be the limitation
itself---the published record of quantized models is outcome-truncated, and we
measured losses roughly four times worse than anything it contains. A risk
estimate built on that record is a floor on risk, not a ceiling, and should say
so.

\bibliographystyle{plainnat}
\bibliography{refs}

\end{document}
